%----------------------------------------------------------------------
\section{Preliminaries}\label{sec:prelim}
%----------------------------------------------------------------------

\subsection{The field $\GF(2)$}

$\GF(2) = \{0, 1\}$ is the field with two elements, with addition
modulo 2. The crucial properties:
\begin{itemize}[nosep]
  \item $1 + 1 = 0$. The additive group is $\mathbb{Z}/2\mathbb{Z}$.
  \item Every element is its own additive inverse: $-x = x$.
  \item There is no negation distinct from identity.
\end{itemize}

The absence of a distinct negation is foundational to what follows.
Boolean complementation requires an involution $x \mapsto \neg x$
with $\neg x \neq x$ for $x \neq 0$. Over $\GF(2)$, this involution
collapses: $\neg 1 = 1 + 1 = 0$, but this is the additive identity,
not a ``complementary truth value'' in the Boolean sense. The
algebraic structure does not inherently support excluded middle.

\subsection{Exterior algebra over $\GF(2)$}\label{sec:exterior}

Let $V = \GF(2)^n$ for some $n$. The exterior algebra
$\bigwedge V = \bigoplus_{k=0}^{n} \bigwedge^k V$ is the graded
algebra generated by $V$ subject to the relation
\begin{equation}\label{eq:self-wedge}
  v \wedge v = 0 \quad \text{for all } v \in V.
\end{equation}

The wedge product is:
\begin{itemize}[nosep]
  \item \textbf{Associative}: $(a \wedge b) \wedge c = a \wedge (b \wedge c)$.
  \item \textbf{Antisymmetric}: $a \wedge b = -b \wedge a$. Over $\GF(2)$,
    since $-1 = 1$, this reduces to $a \wedge b = b \wedge a$.
    The wedge product is \emph{commutative} over $\GF(2)$.
  \item \textbf{Nilpotent on repeated factors}: $v \wedge v = 0$.
\end{itemize}

The identity (\ref{eq:self-wedge}) is the engine of the framework.
It will ensure that self-referential discriminations are algebraically
degenerate.


\subsection{Terminology concordance}
\label{sec:concordance}

The paper draws terminology from several established
frameworks. This concordance specifies the source and
any deviation from standard usage.

\emph{Type theory} (Martin-L\"{o}f 1984; HoTT Book 2013).
\begin{description}[nosep,leftmargin=1.5cm]
  \item[Dependent type] A type whose definition depends on
    a value. Used here for Boolean dimension (\S9):
    $\chi_d(a) = 1$ for all $a$ is a type depending on $d$
    and $S$. Standard usage.
  \item[Type construction / elimination] Introduction and
    use rules for types. Used here (\S5, Remark 5.11):
    the write path constructs types ($\kappa$-evolution
    introduces new discriminators); the read path
    eliminates them (evaluating discriminators on data).
    Analogy, not formal HoTT.
  \item[Univalence] Equivalent types are identical
    (Voevodsky). Used here (Proposition 5.12):
    $\kappa$-extensionality = univalence. The
    operationalist axiom plays the role of the univalence
    axiom. Structural parallel, not a formal embedding.
  \item[Transport] Moving inhabitants between equivalent
    types along a path. Used here (Remark 5.12):
    reclassification of elements under $\kappa$-evolution
    generalizes transport by allowing gap and overlap.
\end{description}

\emph{Logic} (Belnap 1977; Dunn 1976).
\begin{description}[nosep,leftmargin=1.5cm]
  \item[Belnap's four-valued logic / FDE] First Degree
    Entailment: a logic with four truth values (true,
    false, both, neither). Used here (\S9): the subobject
    classifier $\Omega = \GF(2)^2$ recovers exactly these
    four values. Standard identification.
  \item[Paraconsistent] A logic where contradiction does
    not entail everything. Used here (\S4, Remark 4.10):
    the overlap $(1,1)$ tolerates contradiction without
    explosion. Standard usage.
\end{description}

\emph{Category theory} (Mac Lane 1971; Johnstone 2002).
\begin{description}[nosep,leftmargin=1.5cm]
  \item[Topos] A category with finite limits, a subobject
    classifier, and power objects. Used here (Theorem 9.6):
    the fibered category $\int \mathcal{C}$ is claimed to
    be an elementary topos. Standard definition, with the
    caveat that power objects are verified for
    single-discriminator subobjects only (\S B.4).
  \item[Grothendieck construction] Turning a pseudofunctor
    into a fibered category. Used here (Definition 9.4)
    and at the meta-level (Appendix C: the SPPF as a
    fibered category over axiom classes). Standard usage.
  \item[Sheaf / site] A site is a category with a
    Grothendieck topology; a sheaf satisfies the gluing
    condition. Used here (\S9): the site of
    $\kappa$-regimes with six adjunctions as covering
    families. Standard definitions.
  \item[Enriched category] A category whose hom-sets are
    objects of a monoidal category. Used here (Theorem
    9.1): self-enrichment over the category's own monoidal
    structure. Standard usage.
  \item[Adjunction] A pair of functors $F \dashv U$ with
    unit and counit satisfying the zig-zag identities.
    Used here (Theorem 6.11): retract adjunctions from
    direct summand structure. Standard definition;
    the adjunctions are trivial individually.
\end{description}

\emph{Algebraic topology} (Hatcher 2002; May 1999).
\begin{description}[nosep,leftmargin=1.5cm]
  \item[Chain complex] A sequence of abelian groups with
    $\partial \circ \partial = 0$. Used here (\S5):
    the discrimination complex over $\GF(2)$. Standard.
  \item[Homology] $H_k = \ker \partial_k / \mathrm{im}\,
    \partial_{k+1}$. Used here (\S5): residues are
    non-trivial homology classes. Standard.
  \item[$n$-groupoid] A higher groupoid where all
    $k$-morphisms for $k \geq 1$ are invertible. Used
    here (\S5): the symmetric $n$-groupoid from
    $\bigwedge(\GF(2)^n)$. Standard definition.
  \item[Delooping] Viewing a monoidal category as a
    one-object 2-category. Used here (\S8, Theorem 8.3).
    Standard.
\end{description}

\emph{Non-standard usage} (terms specific to this paper).
\begin{description}[nosep,leftmargin=1.5cm]
  \item[Discriminator] A grade-1 element $d \in
    \bigwedge^1 V$ paired with a predicate $d: S \to
    \GF(2)$. After profile linearity, this is a
    \emph{1-form} (covector) over $\GF(2)$: an element
    of $V^* = \operatorname{Hom}(V, \GF(2))$, with the
    evaluation $d(a) = \langle d, a \rangle$ being the
    standard pairing. The term ``discriminator'' adds
    operational semantics to the 1-form: it fires or is
    silent, it is keyed to a datum property, and its
    evaluation is an interrogation act. Note: ``discriminator''
    in universal algebra denotes a ternary operation
    $d(x,y,z)$; the usage here is unrelated.
  \item[$\kappa$-regime] A finite-dimensional subspace of
    $V^*$: the current battery of 1-forms available for
    interrogation. Corresponds to a \emph{context} $\Gamma$
    in type theory ($\Gamma \vdash t : T$) or a
    \emph{stage} in presheaf topos theory (an object of
    the base category). The term ``regime'' emphasizes that
    the subspace governs what can be observed.
  \item[$\kappa$-evolution] Monotonic extension of the
    regime: $\kappa \subseteq \kappa'$. Corresponds to
    \emph{context extension} ($\Gamma \subseteq \Gamma'$)
    in type theory, \emph{refinement} in domain/lattice
    theory, or \emph{base change} in algebraic geometry.
  \item[Residue] $\ker(\chi)$ where $\chi = \tau + \sigma$:
    the elements the current regime cannot resolve. This
    is the \emph{kernel of the XOR map}, a standard
    linear-algebraic object. The paper-specific content
    is the operational interpretation: a non-empty kernel
    triggers $\kappa$-evolution. The term ``residue''
    follows the complex-analysis usage (what remains after
    integration around a singularity): it is what remains
    after discrimination has extracted what it can.
  \item[Coverage profile / $\gamma$-bit] The $\GF(2)^3$
    expansion distinguishing pre-interrogation ($\gamma = 0$)
    from post-interrogation ($\gamma = 1$). Paper-specific;
    closest standard analog is the \emph{support} of a
    function (the set of points where it is non-zero).
  \item[$\GF(2)$ toroid] $\GF(2)^2 = \GF(2) \times \GF(2)$,
    the Klein four-group. The product structure is
    standard; the term ``toroid'' emphasizes that the
    paper systematically shifts between $\GF(2)$ (single
    predicates, the base field) and $\GF(2)^2$ (paired
    predicates, the membership profile). This shift is
    the discrimination process itself: composing two
    $\GF(2)$-valued observations into a $\GF(2)^2$-valued
    profile. Conversely, $\GF(2)$ is the quotient
    $\GF(2)^2 / \GF(2)$ --- the residue from projecting
    the paired profile onto a single axis. The XOR
    signature $\chi = \tau + \sigma$ is this quotient map.
    The ``toroid'' names the product-and-quotient
    structure that the framework traverses, not merely
    the group.
  \item[Cofibration] The outer tensor product of all
    covered domains: each spine claim tensored with each
    appendix verification. An interpretation of the
    tensor (here, the wedge product $\text{Spine} \wedge
    \text{Appendix}$) as a collection of fibers, where
    each fiber is a claim--verification pair. Orphaned
    appendix sections are fibers with no base point;
    uncovered spine theorems are base points with no
    fiber. This is dual to the fibration (Grothendieck
    construction) that the SPPF computes: the SPPF
    fibers formal objects over axiom classes (base $\to$
    total space); the cofibration checks that the total
    space (appendix) covers the base (spine).
\end{description}
